\section{Detailed Audit of the Released CellSAM Checkpoint}\label{detailed-audit-of-the-released-cellsam-checkpoint}

To clarify the actual structure of the released CellSAM checkpoint, we performed both code-level and parameter-level comparisons of its internal branches. The public checkpoint contains three relevant components: \texttt{model}, \texttt{model\_cp}, and \texttt{cellfinder}. For segmentation inference, the released implementation uses the \texttt{model\_cp} branch rather than \texttt{model}. In the same audited path, the mask decoder is called with \texttt{multimask\_output=False}, meaning that the practical segmentation route is the single-mask path rather than the multi-candidate SAM output mode.

Parameter-level comparison further shows that the checkpoint cannot be interpreted simply as a neck-only continuation from \texttt{model} to \texttt{model\_cp}. The branch \texttt{cellfinder.decode\_head.backbone.body} is exactly identical to \texttt{model.image\_encoder} after removing the neck (\texttt{171/171} matching parameter tensors), but entirely different from \texttt{model\_cp.image\_encoder} without the neck (\texttt{0/171}). Likewise, \texttt{model.image\_encoder} and \texttt{model\_cp.image\_encoder} are completely different once the neck is excluded (\texttt{0/171}). At the full-module level, \texttt{model} and \texttt{model\_cp} are globally distinct: the complete image encoder differs in \texttt{177/177} parameter tensors, the neck differs in \texttt{6/6}, the mask decoder differs in \texttt{120/120}, and the prompt encoder differs in \texttt{17/17}.

It is also important to distinguish architecture construction from weight loading. In the released code, \texttt{self.model\ =\ sam\_model\_registry{[}"vit\_b"{]}()} is instantiated first, and \texttt{self.model\_cp} is then created as a deep copy of this freshly constructed model. Therefore, before the released CellSAM checkpoint is loaded, \texttt{model} and \texttt{model\_cp} share the SAM ViT-B architecture and the same fresh initialization state, but they are not automatically the original Meta pretrained SAM checkpoint.

Relative to the original SAM ViT-B checkpoint, the released \texttt{model\_cp} branch preserves only a limited subset of parameters unchanged. The prompt encoder remains fully identical (\texttt{17/17} parameter tensors), while the mask decoder preserves \texttt{30/120} parameter tensors and changes the remaining \texttt{90/120}. The preserved mask-decoder tensors are concentrated in the IoU prediction head, the extra multimask hypernetwork branches, and a few attention \texttt{k\_proj.bias} terms. By contrast, the main single-mask output branch, the output upscaling path, the output tokens, and most transformer blocks are different. Because the released inference path uses \texttt{multimask\_output=False}, this pattern is consistent with further adaptation of the practical single-mask segmentation path, while some unused multi-output components may have remained unchanged.

This appendix supports two bounded claims that are safe for the thesis. First, the released checkpoint semantics must be audited rather than inferred from paper wording alone. Second, the released \texttt{model\_cp} segmentation branch is not the raw SAM mask decoder. What cannot be claimed from the public release alone is the exact hidden training procedure that produced these differences.
